\documentclass[10pt]{article}
\usepackage[a4paper,margin=0.84in]{geometry}
\usepackage[T1]{fontenc}
\usepackage{lmodern}
\usepackage{amsmath,amssymb,amsthm}
\usepackage{graphicx}
\usepackage{microtype}
\usepackage[colorlinks=true,linkcolor=blue,urlcolor=blue,citecolor=blue]{hyperref}
\newtheorem{theorem}{Theorem}
\newtheorem{corollary}{Corollary}
\newtheorem{lemma}{Lemma}
\newcommand{\T}{\mathbb T}
\title{No Infinite-Dimensional Unitary Is Weakly\\
$l$-Sequentially Supercyclic}
\author{Candidate full-result packet\\
\small arXiv:2010.13331, Section 4}
\date{13 August 2026}

\begin{document}
\maketitle

\paragraph{Status.}
Candidate full answer, subject to expert review, to the strengthened
Hilbert-space question in the source.  No unitary on a Hilbert space of
dimension greater than one is weakly $l$-sequentially supercyclic.  It follows
that every weakly $l$-sequentially supercyclic contraction on an
infinite-dimensional Hilbert space is weakly stable.

\section{Source question}
Kubrusly--Vieira ask whether weak supercyclicity implies weak stability for
power-bounded operators, and explicitly sharpen the Hilbert-space version to:
\begin{quote}
``Does there exist a weakly unstable and weakly $l$-sequentially supercyclic
unitary operator?''
\end{quote}

\begin{figure}[ht]
\centering
\includegraphics[width=.95\textwidth]{figures/open_problem_crop.png}
\caption{The broad and strengthened questions in the source PDF, page 5.}
\end{figure}

A vector $y$ is weakly $l$-sequentially supercyclic for $U$ if every $x$ is
the weak limit of a sequence $a_kU^{n_k}y$, where $n_k$ is strictly
increasing.  This is stronger than weak density of the projective orbit.

\section{Proof Intuition}
If a unitary $U$ had such a vector $y$, then $y$ would be cyclic.  In the
cyclic spectral model, $U$ is multiplication by $z$ on $L^2(\mu)$ and $y$ is
the constant function $1$.  A weakly convergent sequence
$a_kz^{n_k}$ is norm bounded, so the scalars $a_k$ are uniformly bounded.
Its weak limit must therefore be essentially bounded.  But an
infinite-dimensional $L^2(\mu)$ always contains unbounded functions.  This
prevents the weak limit set of one projective orbit from being all of the
space.

The distinction is exactly what the Kronecker-set example of
Bayart--Matheron exploits: its projective orbit is weakly dense, but weak
neighbourhoods may require elements of arbitrarily large norm.  Weak density
does not provide a single weakly convergent sequence.

\section{The obstruction}
\begin{lemma}\label{lem:linfty}
Let $\mu$ be a probability measure and $C<\infty$.  The set
\[
 K_C=\{f\in L^2(\mu): |f|\le C\ \mu\text{-a.e.}\}
\]
is weakly closed in $L^2(\mu)$.  If $L^2(\mu)$ is infinite-dimensional, then
$L^\infty(\mu)$ is a proper subset of $L^2(\mu)$.
\end{lemma}

\begin{proof}
$K_C$ is convex.  It is norm closed: from any norm-convergent sequence one
may extract an almost-everywhere convergent subsequence, which preserves the
pointwise bound.  A convex norm-closed subset of a Banach space is weakly
closed.

For the second assertion, if every $L^2$ function were essentially bounded,
the closed graph theorem would make the inclusion
$L^2(\mu)\hookrightarrow L^\infty(\mu)$ bounded.  Thus
$\|f\|_\infty\le M\|f\|_2$.  Applied to indicators, every positive-measure
set would have measure at least $M^{-2}$.  The measure algebra would then
have only finitely many atoms and $L^2(\mu)$ would be finite-dimensional, a
contradiction.
\end{proof}

\begin{theorem}\label{thm:unitary}
No unitary operator on a complex Hilbert space of dimension greater than one
is weakly $l$-sequentially supercyclic.
\end{theorem}

\begin{proof}
First suppose the Hilbert space $H$ is infinite-dimensional.  If $y$ were a
weakly $l$-sequentially supercyclic vector, the norm-closed linear span of
$\{U^ny:n\ge0\}$ would be all of $H$: that closed subspace is weakly closed
and contains every weak sequential limit of the projective orbit.
Consequently $y$ is cyclic.

The cyclic spectral theorem gives a probability measure $\mu$ on $\T$ and a
unitary map $W:H\to L^2(\mu)$ such that
\[
 Wy=1,\qquad WUW^{-1}=M_z.
\]
Choose $g\in L^2(\mu)\setminus L^\infty(\mu)$ using Lemma
\ref{lem:linfty}.  Weak $l$-sequential supercyclicity would give scalars
$a_k$ and increasing integers $n_k$ such that
\[
 a_kz^{n_k}\rightharpoonup g\quad\text{in }L^2(\mu).
\]
Every weakly convergent sequence in a Banach space is norm bounded.  Since
$\|a_kz^{n_k}\|_2=|a_k|$, there is a finite $C$ with $|a_k|\le C$ for all
$k$.  Hence every term belongs to $K_C$, and weak closedness of $K_C$ forces
$g\in K_C\subset L^\infty(\mu)$, a contradiction.

If $2\le\dim H<\infty$, diagonalize $U$.  Relative to an eigenbasis, the
absolute values of the coordinates of $U^ny/\|y\|$ are independent of $n$.
Thus the normalized projective orbit cannot be dense in the unit sphere, so
it cannot be weakly $l$-sequentially supercyclic.  (In finite dimensions weak
and norm convergence agree.)
\end{proof}

\begin{corollary}\label{cor:contraction}
Every weakly $l$-sequentially supercyclic contraction on an
infinite-dimensional Hilbert space is weakly stable.
\end{corollary}

\begin{proof}
Use the Sz.-Nagy--Foia\c{s}--Langer decomposition quoted in Remark 4.1 of the
source:
\[
 H=H_c\oplus H_u,\qquad T=C\oplus U,
\]
where $C$ is completely nonunitary and weakly stable, while $U$ is unitary.
Let $y=y_c\oplus y_u$ be weakly $l$-sequentially supercyclic.  Projection of
approximating sequences onto $H_u$ shows that, if $\dim H_u>1$, then $y_u$
is weakly $l$-sequentially supercyclic for $U$, contradicting Theorem
\ref{thm:unitary}.

If $\dim H_u=1$, write $U=\lambda I$ and choose nonzero $x\in H_c$.  To
converge weakly to $x\oplus1$, a sequence
\[
 a_k(C^{n_k}y_c\oplus \lambda^{n_k}y_u)
\]
must have $a_k\lambda^{n_k}y_u\to1$, so $(a_k)$ is bounded.  Weak stability
of $C$ then gives $a_kC^{n_k}y_c\rightharpoonup0$, contradicting convergence
to $x$.  Hence $H_u=\{0\}$ and $T=C$ is weakly stable.
\end{proof}

\section{Verification, literature audit, and limitations}
The proof is analytic and uses no numerical computation.  Its two audit
points are: (i) weak convergence of a sequence implies norm boundedness; and
(ii) the uniformly $L^\infty$-bounded ball is weakly closed in $L^2$.

A bounded search through 13 August 2026 checked the run indexes, the exact
question, the 2024 survey arXiv:2405.02752, the 2023 paper
arXiv:2306.08197, the Bayart--Matheron construction, and searches combining
``weakly $l$-sequentially supercyclic'', ``unitary'', and ``weak stability''.
No version of Theorem \ref{thm:unitary} was found.  Several later papers
instead describe the Bayart--Matheron unitary as weakly
$l$-sequentially supercyclic.  The displayed argument shows why weak density
of that orbit cannot imply this stronger direct-sequence property.  This is
not a claim of priority and the conflict makes independent review essential.

The source's broader question for merely weak supercyclic operators is not
answered positively by this theorem.  In fact, the same Kronecker-set
construction gives a literature-implied negative answer to that broader
question; it is recorded in a companion packet.

\paragraph{Human-review recommendation.}
Verify the cyclic spectral representation with $y\mapsto1$, the weak
closedness of $K_C$, and the one-dimensional unitary-summand argument in
Corollary \ref{cor:contraction}.  Then compare the direct-sequence definition
in the source with the weak-density argument in Bayart--Matheron Example 3.6.

\begin{thebibliography}{9}
\bibitem{KV}
C. S. Kubrusly and P. C. M. Vieira,
\emph{Boundedly Spaced Subsequences and Weak Dynamics},
Journal of Function Spaces (2018), Article ID 4732836;
\href{https://arxiv.org/abs/2010.13331}{arXiv:2010.13331}.

\bibitem{BM}
F. Bayart and E. Matheron,
\emph{Hyponormal operators, weighted shifts and weak forms of
supercyclicity}, Proc. Edinburgh Math. Soc. 49 (2006), 1--15,
doi:10.1017/S0013091504000975.
\end{thebibliography}
\end{document}
